\section{Proof of the main theorem}\label{sec:conclusion}

\begin{proof}[Proof of \cref{thm:main}]
The global Kneser--Tits theorem handles the isotropic case, and the known
inner-type-$A$ theorem handles anisotropic inner forms.  As explained in
\cref{sec:arithmetic}, Harder's theorem leaves only the anisotropic
unitary groups $S=\SU(D,*)$ of degree $n\geq3$.

Let $N\triangleleft S(k)$ be noncentral.  By
\cref{arith:finite-index}, the group $S(k)/N$ is finite.  Choose an integer
$e\geq1$ divisible by its exponent.  With the notation
\eqref{arith:power-notation}, we have $R\subset N$, while
\cref{ind:main} gives
\[
 R=V_0=\delta_A^{-1}(P_0).
\]
Thus $N$ is a union of cosets of the open subgroup
$\delta_A^{-1}(P_0)$.  More explicitly, density of $\delta_A(S(k))$
induces an isomorphism of finite groups
\[
 S(k)/R\ \simeq\ H_A/P_0.
\]
Let $W$ be the inverse image in $H_A$ of the normal subgroup corresponding
to $N/R$.  Then $W$ is open and normal, and
$N=\delta_A^{-1}(W)$.  Density also shows that
$W=\overline{\delta_A(N)}$, so this local subgroup is uniquely determined
by $N$.  If $A$ is empty, $H_A$ is trivial and the same argument gives
$N=S(k)$.
\end{proof}
